\section{Italian Units And Allied Progress}\label{sec:10.0}

\ruleslabel{General Rule}

\begin{generalrule}
Certain events outside the scope of the Balkans campaign trigger pullbacks, withdrawals, and surrender of Italian units, as well as the entry of special German reinforcements. These events are measured on the Allied Progress Track.
\end{generalrule}

\ruleslabel{Cases}

\subsection{Allied Progress}\label{sec:10.1}

Starting with Game-Turn 6, the Yugoslav Player rolls a single die during the Allied Progress Phase and applies the result to the Allied Progress Track (10.14). There are seven events displayed on this Track, which, as they occur, may trigger events affecting Italian units.

\case{10.11} At the beginning of Game-Turn 6, the Allied Progress marker is placed in the upper-half of Box \#1 (Alam Halfa) on the Allied Progress Track. When the Yugoslav Player rolls the die during the Allied Progress Phase, a roll of 1 through 5 indicates that the marker is advanced to the upper half of the next box, where it remains until the next Allied Progress Phase. A roll of 6 indicates that the marker is not advanced but is simply shifted to the lower-half of the box it currently occupies, indicating that the Allied campaign has been temporarily stalled.

\case{10.12} If the Allied Progress marker currently occupies the lower-half of a box, one is added to the Yugoslav Player’s die roll during the Allied Progress Phase. The maximum addition to this die roll is One, no matter how long the marker remains in the lower-half of the box. No modification to the die roll is made when the marker occupies the upper-half of the box.

\case{10.13} If the Allied Progress marker currently occupies the lower-half of a box and a 6 or 7 is rolled by the Yugoslav Player during the Allied Progress Phase, the marker is not moved at all; it remains in the lower-half of the box it currently occupies (see 10.12). Whenever the marker is advanced, it is always placed in the upper-half of the next box.

\case{10.14} Allied Progress Track (see mapsheet)

\subsection{Italian Pullbacks And Withdrawals}\label{sec:10.2}

\case{10.21} When the Allied Progress marker enters Box \#3 on the Allied Progress Track, the Axis Player is immediately subject to an Italian Pullback for the remainder of the game. When an Italian Pullback is in effect, Italian units in Croatia may not enter the Zone Display for any reason, and may only occupy Objective Displays labeled ‘‘Italian Pullback’’ in the Axis box. Italian units currently occupying a prohibited location at the moment an Italian Pullback comes into effect must move out in the ensuing Axis Movement Phase. Note: Once an Italian Pullback has occurred, Italian units are not permitted to participate in Anti-Guerrilla Operations in Bosnia or Serbia (see 6.64B).

\case{10.22} When the Allied Progress marker enters Box \#6 on the Allied Progress Track, the Axis Player is immediately subject to an Italian Withdrawal. In this event, the Axis Player immediately rolls a single die, the resulting number indicating the number of Italian units (Axis Player’s choice) which must be immediately and permanently removed from the map. This die roll takes place only once per game. In the ensuing Axis Reinforcement Phase, German reinforcements due as a result of Italian Withdrawal (see 13.91) become available to the Axis Player.

\case{10.23} No Italian unit may participate in an Anti-Guerrilla Operation once an Italian Withdrawal takes place.

\subsection{Italian Surrender And Operation Konstantin}\label{sec:10.3}

When the Allied Progress marker reaches Box \#7 on the Allied Progress Track, Italian Surrender takes place. As a result, all Italian units on the map either disband or defect. This determination is made during the Italian Surrender Phase. This Phase occurs only once per game in the Game-Turn in which Box \#7 is reached. After Italian Surrender, the Allied Progress Phase is omitted for the rest of the game. See 7.23B.

\case{10.31} The Axis Player receives reinforcements in the Axis Reinforcement Phase following the Allied Progress Phase in which Italian Surrender occurs (see 13.91).

\case{10.32} In the Italian Surrender Phase, the following procedure is performed:

\textbf{A.} All Italian units on the map either disband (see 10.33), defect (see 10.34), or remain in place (see 10.35).

\textbf{B.} All German units perform Operation Konstantin (see 10.36).

\textbf{C.} All Partisan units respond to Operation Konstantin (see 10.36).

\textbf{D.} All Italian units that remained in place according to the procedure in Step A either disband, defect, or continue to remain in place.

\textbf{E.} All remaining Italian units on the map are permanently removed.

\case{10.33} During Steps A and D of Case 10.32, a given Italian unit disbands if it occupies an Axis box, triangle, or circle containing an equal amount or more non-Italian Axis Strength Points than Partisan Strength Points in a corresponding box, triangle, or circle. A disbanded unit is immediately removed from the map.

\case{10.34} During Steps A and D of Case 10.32, a given Italian unit defects if
it occupies an Axis box, triangle, or circle with fewer non-Italian Axis Strength
Points than Partisan Strength Points in a corresponding box, triangle, or circle.
If an Italian unit defects, it is immediately removed from the map and replaced
with a Partisan brigade (if available), which is placed in the Yugoslav box,
triangle, or circle corresponding to the one formerly occupied by the Italian
unit.

\case{10.35} During Steps A and D of Case 10.32, a given Italian unit remains in place if it occupies an Axis box, triangle, or circle without any other non-Italian Axis units and no Partisan units in a corresponding location.

\case{10.36} When Operation Konstantin comes into effect (see 10.32B), the Axis Player is permitted a bonus Movement Phase for German units only. This bonus Phase is identical to a normal Movement Phase. As soon as all German units have completed this movement, the Yugoslav Player is permitted an identical bonus Movement Phase for all Partisan units. This is also considered a normal Movement Phase, except the Yugoslav Player is not subject to the restrictions of Case 6.42 (i.e., a Partisan unit may move into a Zone and enter one of its Objective Displays immediately). Bonus Movement Phases are in addition to normal movement permitted during Axis and Yugoslav Movement Phases.
